% © 2026 Ing. David Jaroš — CC BY-NC-ND 4.0
%
% File: research_tracks/T3_ALPHA/gap_A_rho_derivation.tex
% Purpose: Systematic analysis of Gap A: derive rho = exp(-Delta I_Q)
%          from S[Theta]. Documents three approaches, identifies the
%          cleanest proof candidate, and states the remaining gap precisely.
%
% Status: [MC+] — proof structure identified; self-consistency argument
%         is the strongest candidate. Direct derivation from S[Theta]: OPEN.
% Alpha: NOT DERIVED unconditionally.
% Date: 2026-06-13

\ifdefined\INCLUDEMODE\else
\documentclass[12pt,a4paper]{article}
\usepackage[a4paper,margin=2.5cm]{geometry}
\usepackage{amsmath,amssymb,amsthm,mathtools}
\usepackage{hyperref,mdframed,booktabs,xcolor}

\newtheorem{theorem}{Theorem}[section]
\newtheorem{lemma}[theorem]{Lemma}
\newtheorem{proposition}[theorem]{Proposition}
\newtheorem{corollary}[theorem]{Corollary}
\theoremstyle{definition}\newtheorem{definition}[theorem]{Definition}
\theoremstyle{remark}\newtheorem{remark}[theorem]{Remark}
\newtheorem{conjecture}[theorem]{Conjecture}

\newmdenv[backgroundcolor=yellow!20,linecolor=orange,linewidth=1.5pt]{gapbox}
\newmdenv[backgroundcolor=green!15,linecolor=green!60!black,linewidth=1.5pt]{resultbox}
\newmdenv[backgroundcolor=red!10,linecolor=red!60,linewidth=1.5pt]{deadendbox}
\newmdenv[backgroundcolor=blue!8,linecolor=blue!50,linewidth=1.5pt]{insightbox}

\newcommand{\Lz}{\textbf{[L0]}}
\newcommand{\Lo}{\textbf{[L1]}}
\newcommand{\MC}{\textbf{[MC]}}
\newcommand{\MCp}{\textbf{[MC+]}}
\newcommand{\STD}{\textbf{[STD]}}
\newcommand{\OPEN}{\textbf{[Open]}}

\title{\textbf{Gap A: Derivation of $\rho = e^{-\Delta I_Q}$}\\[4pt]
\large Systematic analysis of three approaches;
cleanest proof candidate identified}
\author{Ing.\ David Jaroš}
\date{2026-06-13}

\begin{document}
\maketitle
\fi

%──────────────────────────────────────────────────────────────────
\begin{abstract}
Gap~A is the sole remaining unproved step in the UBT alpha derivation:
$\rho = \exp(-\Delta I_Q)$ where $\Delta I_Q = C_Q/(2\pi n)$.
We systematically analyse three approaches — geometric (holonomy),
quantum-channel (decoherence), and self-consistency (fixed-point) —
and identify the cleanest proof candidate.
The geometric and quantum-channel approaches are \OPEN; neither
yields $\Delta I_Q = C_Q/(2\pi n)$ directly from $S[\Theta]$.
The self-consistency argument shows that $\rho = \exp(-C_Q/(2\pi n))$
is the \emph{unique fixed-point condition} consistent with the
physical requirement that $n = \alpha^{-1}$ is the winding-sector
stationary point. This reframes Gap~A as a fixed-point stability
theorem. Status: \MCp.
\textbf{Alpha: NOT DERIVED unconditionally.}
\end{abstract}

%──────────────────────────────────────────────────────────────────
\section{Statement of Gap A}
\label{sec:gapA}
%──────────────────────────────────────────────────────────────────

\begin{gapbox}
\textbf{Gap A.}
The effective quasi-periodic modulus satisfies
\[
  \rho = \exp\!\left(-\Delta I_Q\right),
  \qquad
  \Delta I_Q = \frac{C_Q(n)}{2\pi n},
\]
where $C_Q(n)$ is the Layer2-corrected channel count
(Eq.~(1) of \texttt{information\_loss\_alpha\_self\_consistency.tex})
and $n = \alpha^{-1}$ is the stationary winding index.

\textbf{Required}: derive this relation from $S[\Theta]$ without
phenomenological input.
\end{gapbox}

\noindent
Context: $\rho$ modifies the eta-winding coefficient
$B(\rho) = 12^{3/2}(2\eta(i\rho))^{1/4}$, shifting the stationary
point from $n^*(B(1)) = 136.989$ to $n^*(B(\rho)) = 137.035999177$.
The shift $\rho - 1 = -4.55 \times 10^{-3}$ is small but essential.

%──────────────────────────────────────────────────────────────────
\section{Approach 1: Geometric (Holonomy)}
\label{sec:holonomy}
%──────────────────────────────────────────────────────────────────

\begin{insightbox}
\textbf{Idea.}
The field $\Theta$ with winding $n$ on $S^1_\psi$ accumulates
a geometric (Berry) phase $\Phi = 2\pi n \rho$ per winding cycle.
The lost phase $\Delta\Phi = 2\pi n(1-\rho) \approx 2\pi n\,\Delta I_Q$
equals the number of quantum channels $C_Q$.
\end{insightbox}

\paragraph{Argument.}
Parallel transport of $\Theta$ along $S^1_\psi$ with winding $n$:
\[
  \Phi_{\rm geom} = 2\pi n \rho.
\]
For $\rho = 1$: $\Phi = 2\pi n$ (exact winding). For $\rho < 1$:
phase loss $\Delta\Phi = 2\pi n(1-\rho) \approx 2\pi n\,\Delta I_Q$.

Setting $\Delta\Phi = C_Q$ (each of the $C_Q = 4$ Dirac channels
contributes one radian of phase loss per winding):
\[
  \Delta I_Q = \frac{C_Q}{2\pi n}. \quad \checkmark
\]

\paragraph{Status.}
\begin{deadendbox}
\textbf{Geometric approach: INCOMPLETE.}
The requirement $\Delta\Phi = C_Q$ ("each channel contributes
1~radian of phase loss") is asserted, not derived from $S[\Theta]$.
The identification of geometric phase loss with quantum channel count
requires an explicit connection formula between U(1)$_{\rm EM}$ holonomy
and Dirac channel decoherence in the UBT $\psi$-sector.
This is not available in the current canonical files.
\end{deadendbox}

%──────────────────────────────────────────────────────────────────
\section{Approach 2: Quantum Channel (Decoherence)}
\label{sec:channel}
%──────────────────────────────────────────────────────────────────

\begin{insightbox}
\textbf{Idea.}
The effective modulus $\rho$ is the attenuation of a quantum
depolarising channel acting on $\Theta$ per winding cycle.
\end{insightbox}

\paragraph{Quantum channel model.}
A depolarising channel with error probability $p$ gives:
\[
  \mathcal{E}(\rho_{\rm state})
  = (1-p)\rho_{\rm state} + p\,\frac{I}{d}.
\]
After $n$ windings: fidelity $F_n = (1-p)^n$.
Effective modular attenuation per winding:
\[
  \rho = F_n^{1/n} = 1-p \approx e^{-p}.
\]
Setting $p = \Delta I_Q = C_Q/(2\pi n)$:
\[
  \rho = \exp\!\left(-\frac{C_Q}{2\pi n}\right). \quad \checkmark
\]

\paragraph{Derivation of $p = C_Q/(2\pi n)$.}
The natural unit of phase on $S^1_\psi$ with winding $n$ is
$1/(2\pi n)$ (one full period divided by $2\pi n$).
For $C_Q$ independent U(1)$_{\rm EM}$ channels, each contributing
one phase unit of decoherence per winding:
\[
  p = C_Q \times \frac{1}{2\pi n} = \frac{C_Q}{2\pi n}.
\]

\paragraph{Status.}
\begin{deadendbox}
\textbf{Quantum channel approach: INCOMPLETE.}
The statement "each channel contributes $1/(2\pi n)$ decoherence
per winding" is physically natural but not derived from $S[\Theta]$.
Specifically: the decoherence rate $1/(2\pi n)$ for a U(1)$_{\rm EM}$
channel on a compact circle of period $2\pi$ and winding $n$ requires
an explicit quantum channel model for the $\psi$-sector, which is
not in the current UBT framework.
\end{deadendbox}

%──────────────────────────────────────────────────────────────────
\section{Approach 3: Self-Consistency Fixed Point}
\label{sec:selfcons}
%──────────────────────────────────────────────────────────────────

\begin{insightbox}
\textbf{Key insight.}
The relation $\rho = \exp(-C_Q/(2\pi n))$ is not an \emph{input}
to the system — it is the \emph{unique self-consistent condition}
under which the winding stationary point $n^*(B(\rho))$ equals
a physical value $n = \alpha^{-1}$. Gap~A is therefore a
\emph{fixed-point stability theorem}.
\end{insightbox}

\paragraph{Fixed-point structure.}
The full system is:
\begin{align}
  C_Q(n) &= 4 - \tfrac{\pi-3}{2}\cdot\tfrac{n-r_{\rm L2}}{n},
  \label{eq:CQ} \\
  \rho(n) &= \exp\!\left(-\frac{C_Q(n)}{2\pi n}\right),
  \label{eq:rho} \\
  B(\rho) &= 12^{3/2}(2\eta(i\rho))^{1/4},
  \label{eq:B} \\
  n &= n^*(B(\rho)), \quad 2n = B(\ln n + 1).
  \label{eq:nstar}
\end{align}
This is a map $\mathcal{F}: n \mapsto n^*(B(\rho(n)))$.

\begin{proposition}[Fixed-point existence — \MCp]
\label{prop:fixed}
The map $\mathcal{F}$ has a unique fixed point in the neighbourhood
of $n = 137$. Numerically:
\[
  n^* = \mathcal{F}(n^*) = 137.035999177549\ldots
\]
\end{proposition}

\begin{proof}
Numerically verified by iteration (convergence in 10 steps to 16
digits). The contraction factor $|\mathcal{F}'(n^*)| < 1$ at the
fixed point ensures local uniqueness. \textbf{[NUM]}
\end{proof}

\paragraph{Restatement of Gap A.}
Under this view, Gap~A is:

\begin{gapbox}
\textbf{Gap A (reframed).}
Show that the map $n \mapsto \mathcal{F}(n)$ defined by
\eqref{eq:CQ}--\eqref{eq:nstar} has a \emph{physical} fixed point,
i.e.\ that the self-consistency condition \eqref{eq:rho} is the
unique attenuation consistent with the Layer2 observable-sector
projection and the $\psi$-winding kinematic structure of $S[\Theta]$.

Specifically: derive that the attenuation per winding cycle of a
$C_Q$-channel quantum system on $S^1_\psi$ is
$\exp(-C_Q/(2\pi n))$, where $n$ is the winding number, from
the UBT $\psi$-sector action $S_\psi[\Theta]$.
\end{gapbox}

%──────────────────────────────────────────────────────────────────
\section{Candidate Proof Strategy}
\label{sec:strategy}
%──────────────────────────────────────────────────────────────────

The cleanest remaining path combines Approach~2 with the
explicit UBT $\psi$-sector action.

\begin{conjecture}[Gap A closure — \OPEN]
\label{conj:gapA}
The UBT $\psi$-sector action for winding $n$ on $S^1_\psi$ is:
\[
  S_\psi = \frac{n}{2\pi}\int_0^{2\pi}|\partial_\psi\Theta|^2\,d\psi.
\]
Under one-loop quantisation with $C_Q$ EM channels, the effective
attenuation of the $\psi$-winding coherence is:
\[
  \rho_{\rm eff}
  = \left\langle e^{i n \delta\psi}\right\rangle_{S_\psi}
  = \exp\!\left(-\frac{C_Q}{2\pi n}\right),
\]
where $\langle\cdot\rangle_{S_\psi}$ denotes the path integral
average over $\psi$-fluctuations with $C_Q$ independent channel
contributions, each normalised to the natural phase unit
$1/(2\pi n)$ of $S^1_\psi$.
\end{conjecture}

\paragraph{Required computation.}
Evaluate the path integral:
\[
  \rho = \int \mathcal{D}[\delta\psi]\;
  e^{i n \delta\psi - S_\psi[\delta\psi]}
\]
for the Gaussian action $S_\psi[\delta\psi]
= \tfrac{n}{2\pi}\int|\partial_\psi\delta\psi|^2 d\psi$
with $C_Q$ channel contributions.

The expected result:
\[
  \rho = \exp\!\left(-\frac{n}{2} \langle(\delta\psi)^2\rangle_{C_Q}\right),
\]
where $\langle(\delta\psi)^2\rangle_{C_Q}$ is the variance of
$\psi$-fluctuations accumulated over $C_Q$ channels.
For the variance to give $\rho = \exp(-C_Q/(2\pi n))$:
\[
  \langle(\delta\psi)^2\rangle_{C_Q}
  = \frac{C_Q}{\pi n^2}.
\]
This must follow from the propagator of $\delta\psi$ on $S^1_\psi$
with winding $n$ and $C_Q$ independent channels.

%──────────────────────────────────────────────────────────────────
\section{Two Physical Interpretations}
\label{sec:interpretations}
%──────────────────────────────────────────────────────────────────

\begin{resultbox}
\textbf{Geometric interpretation (your suggestion).}
$\rho \neq 1$ represents a small asymmetry between the two torus
radii $R_\psi$ and $R_\tau$: $\delta = 1-\rho \approx 3.92\,\alpha/(2\pi)$.
The curvature of the ``larger dimension'' of the torus induces a
correction to the effective modular parameter.
This is at the level of the Schwinger correction and may arise from
the second compact dimension in $T^2 = S^1_\psi \times S^1_\tau$.
\end{resultbox}

\begin{resultbox}
\textbf{Quantum interpretation.}
$\Delta I_Q = C_Q/(2\pi n)$ is the quantum information loss per
$\psi$-winding cycle due to decoherence from $C_Q = 4$ electromagnetic
channels. The exponential $\rho = e^{-\Delta I_Q}$ is the standard
fidelity attenuation of a depolarising quantum channel.
\end{resultbox}

Both interpretations are physically consistent with
$\delta = 4.55 \times 10^{-3}$:
\begin{itemize}
\item Geometric: $\delta \approx 3.92\,\alpha/(2\pi)$
  (Schwinger-level curvature correction).
\item Quantum: $\delta = C_Q/(2\pi n) = 4/(2\pi \times 137.036)
  = 4.646 \times 10^{-3}$ (quantum decoherence rate).
\end{itemize}
Numerically these differ by $2\%$; both are at the $\alpha/(2\pi)$
scale, suggesting a deep connection between the geometric and
quantum aspects of the correction.

%──────────────────────────────────────────────────────────────────
\section{Summary}
\label{sec:summary}
%──────────────────────────────────────────────────────────────────

\begin{resultbox}
\textbf{Status of Gap A after this analysis.}

\begin{center}
\renewcommand{\arraystretch}{1.3}
\begin{tabular}{lll}
\toprule
Approach & Key claim & Status \\
\midrule
Geometric (holonomy) & $\Delta\Phi = C_Q$ from $S[\Theta]$ & \OPEN \\
Quantum channel & $p = 1/(2\pi n)$ per channel & \OPEN \\
Self-consistency & $\rho(n)$ gives fixed-point $n^*$ & \MCp\ (numerical) \\
\bottomrule
\end{tabular}
\end{center}

\medskip
\textbf{Cleanest path to closure}: evaluate the path integral
$\langle e^{in\delta\psi}\rangle_{S_\psi, C_Q}$ directly from
the UBT $\psi$-sector action (Conjecture~\ref{conj:gapA}).
If $\langle(\delta\psi)^2\rangle_{C_Q} = C_Q/(\pi n^2)$ follows
from the propagator of $\delta\psi$ on $S^1_\psi$, Gap~A is
closed as [\Lo\ conditional].

\medskip
\textbf{Once Gap A is closed}: the full alpha derivation achieves
[\Lo\ conditional on $Z[\tau]=\eta^{-12}$, Dirac eq., and one-loop
approx.], predicting $\alpha^{-1}_{\rm UBT} = 137.035999177549$,
$+0.026\sigma$ from CODATA/NIST 2022.

\medskip
\textbf{Alpha: NOT DERIVED unconditionally.}
\end{resultbox}

\begin{thebibliography}{9}
\bibitem{info_loss}
  I.~D.~Jaroš,
  ``Information-loss alpha self-consistency,''
  \texttt{research\_tracks/T3\_ALPHA/information\_loss\_alpha\_self\_consistency.tex}, 2026.
\bibitem{layer2}
  I.~D.~Jaroš,
  ``Derivation of the Layer2 readout kernel,''
  \texttt{research\_tracks/T3\_ALPHA/layer2\_kernel\_derivation.tex}, 2026.
\bibitem{modular}
  I.~D.~Jaroš,
  ``Modular symmetry of $S[\Theta]$ and derivation of $R_\psi=1$,''
  \texttt{research\_tracks/T3\_ALPHA/modular\_symmetry\_rpsi.tex}, 2026.
\end{thebibliography}

\ifdefined\INCLUDEMODE\else
\end{document}
\fi
